% =====================================================================
%  1.1 Potenser — ekstra utfordring   (MØNSTERSEKSJON)
% =====================================================================

\begin{infoboks}{Velg metode selv}
Arket bygger på 1.1 Potenser. Vis regning og begrunnelse. Bruk eksakte svar.
Kontroller svaret ved å sette inn tall eller regne på en annen måte.
\end{infoboks}

\nivaahode{2}

\oppg{U1}{2}{\uh}
Forenkle. Skriv svaret uten negative eksponenter.
\[ \frac{(2a^2b^{-1})^3 \cdot (ab)^2}{4a^7b^{-2}} \]

\oppg{U2}{2}{\uh}
Skriv alle grunntallene som produkter av 2 og 3. Regn deretter ut.
\[ \frac{2^5 \cdot 6^3}{3^2 \cdot 12^2} \]

\oppg{U3}{2}{\uh}
Løs likningene. Skriv begge sider med samme grunntall først.
\begin{tasks}(3)
  \task $2^x = \dfrac{1}{32}$
  \task $9^x = 27$
  \task $4^{x+1} = 8^x$
\end{tasks}

\oppg{U4}{2}{\prog}
Se på programmet.
\begin{kode}
n = 0
while 2**n < 10**6:
    n = n + 1
print(n)
\end{kode}
\begin{tasks}(1)
  \task Hva skriver programmet ut? Forklar uten å kjøre det.
  \task Endre programmet slik at det finner den minste $n$ der $3^n > 10^9$. Hva blir svaret?
\end{tasks}

\nivaahode{3}

\oppg{U5}{3}{\uh}
Hvilket tall er størst, $2^{100}$ eller $3^{60}$? Begrunn uten kalkulator.
\hint{Skriv begge som en potens med eksponent 20.}

\oppg{U6}{3}{\uh}
Hvor mange sifre har tallet $4^6 \cdot 5^{10}$? Begrunn uten å regne ut hele tallet.

\oppg{U7}{3}{\uh}
Ola påstår at $(a^m)^n = a^{(m^n)}$.
\begin{tasks}(1)
  \task Finn et eksempel som viser at påstanden er feil.
  \task Finn alle par av positive heltall $m$ og $n$ der påstanden likevel stemmer. Begrunn at du har funnet alle.
\end{tasks}

\begin{oppsummering}
Skriv om til samme grunntall. Bruk én regel per steg.
En negativ eksponent gir en brøk, ikke et negativt tall.
Kontroller ved å sette inn et tall.
\end{oppsummering}
